\chapter{Index des méthodes}

\noindent\textit{Les 112 méthodes du fascicule, regroupées par chapitre, pour retrouver
rapidement « comment faire » sans connaître le numéro du chapitre.}
\vspace{8pt}

\begin{multicols}{2}
\footnotesize

\textbf{\normalsize\color{onbo}Chapitre 1 --- Nombres rationnels et opérations}
\begin{enumerate}[label=\arabic*.,itemsep=1pt,topsep=2pt,leftmargin=1.6em]
\item Simplifier une fraction jusqu'à sa forme irréductible
\item Comparer deux fractions
\item Additionner ou soustraire des fractions de dénominateurs différents
\item Multiplier et diviser des fractions
\item Appliquer correctement les priorités opératoires
\item Utiliser la valeur absolue pour comparer des rationnels
\item Résoudre un problème concret avec des rationnels
\end{enumerate}
\vspace{4pt}

\textbf{\normalsize\color{onbo}Chapitre 2 --- Puissances}
\begin{enumerate}[label=\arabic*.,itemsep=1pt,topsep=2pt,leftmargin=1.6em]
\item Appliquer les règles de calcul sur des puissances de même base
\item Calculer une puissance d'exposant négatif
\item Écrire un nombre en notation scientifique
\item Passer de la notation scientifique à l'écriture décimale
\item Calculer un produit ou un quotient en notation scientifique
\item Encadrer / estimer un ordre de grandeur
\item Simplifier une expression combinant plusieurs puissances de même base
\end{enumerate}
\vspace{4pt}

\textbf{\normalsize\color{onbo}Chapitre 3 --- Racines carrées}
\begin{enumerate}[label=\arabic*.,itemsep=1pt,topsep=2pt,leftmargin=1.6em]
\item Simplifier une racine carrée
\item Additionner ou soustraire des racines carrées
\item Multiplier ou diviser des racines carrées
\item Rendre rationnel un dénominateur
\item Comparer deux nombres contenant des racines carrées
\item Résoudre une équation du type $x^{2}=a$
\item Résoudre un problème géométrique avec des racines carrées
\end{enumerate}
\vspace{4pt}

\textbf{\normalsize\color{onbo}Chapitre 4 --- Calcul littéral}
\begin{enumerate}[label=\arabic*.,itemsep=1pt,topsep=2pt,leftmargin=1.6em]
\item Développer avec la distributivité simple
\item Développer avec la double distributivité
\item Développer avec une identité remarquable
\item Factoriser par un facteur commun évident
\item Factoriser un facteur commun "caché" dans une somme de produits
\item Factoriser avec une identité remarquable
\item Utiliser une identité remarquable pour un calcul astucieux
\end{enumerate}
\vspace{4pt}

\textbf{\normalsize\color{onbo}Chapitre 5 --- Équations et inéquations du premier degré}
\begin{enumerate}[label=\arabic*.,itemsep=1pt,topsep=2pt,leftmargin=1.6em]
\item Résoudre une équation $ax+b=cx+d$
\item Résoudre une équation avec parenthèses
\item Résoudre une équation produit nul
\item Factoriser puis résoudre une équation produit nul
\item Résoudre une inéquation
\item Représenter une inéquation sur une droite graduée
\item Mettre un problème en équation ou en inéquation
\end{enumerate}
\vspace{4pt}

\textbf{\normalsize\color{onbo}Chapitre 6 --- Systèmes de deux équations à deux inconnues}
\begin{enumerate}[label=\arabic*.,itemsep=1pt,topsep=2pt,leftmargin=1.6em]
\item Résoudre un système par substitution
\item Résoudre un système par combinaison linéaire
\item Choisir la méthode la plus adaptée
\item Résoudre un système où les deux équations doivent être multipliées
\item Interpréter graphiquement la solution d'un système
\item Mettre un problème concret en système
\item Vérifier une solution de système
\end{enumerate}
\vspace{4pt}

\textbf{\normalsize\color{onbo}Chapitre 7 --- Arithmétique : PGCD, PPCM, nombres premiers}
\begin{enumerate}[label=\arabic*.,itemsep=1pt,topsep=2pt,leftmargin=1.6em]
\item Décomposer un nombre en produit de facteurs premiers
\item Calculer un PGCD par l'algorithme d'Euclide
\item Calculer un PPCM
\item Rendre une fraction irréductible
\item Reconnaître si deux nombres sont premiers entre eux
\item Résoudre un problème de partage (PGCD)
\item Résoudre un problème de synchronisation (PPCM)
\end{enumerate}
\vspace{4pt}

\textbf{\normalsize\color{onbo}Chapitre 8 --- Proportionnalité, pourcentages, échelles}
\begin{enumerate}[label=\arabic*.,itemsep=1pt,topsep=2pt,leftmargin=1.6em]
\item Compléter un tableau de proportionnalité
\item Calculer un pourcentage d'une quantité
\item Calculer un taux de variation
\item Appliquer une évolution en pourcentage
\item Enchaîner deux évolutions successives
\item Utiliser une échelle
\item Calculer une vitesse moyenne, une distance ou une durée
\end{enumerate}
\vspace{4pt}

\textbf{\normalsize\color{onbo}Chapitre 9 --- Fonctions linéaires et fonctions affines}
\begin{enumerate}[label=\arabic*.,itemsep=1pt,topsep=2pt,leftmargin=1.6em]
\item Calculer l'image d'un nombre par une fonction
\item Déterminer un antécédent
\item Reconnaître une fonction linéaire ou affine, identifier $a$ et $b$
\item Tracer la représentation graphique d'une fonction affine
\item Déterminer une fonction affine par deux points
\item Déterminer une fonction affine par un point et le coefficient directeur
\item Comparer deux tarifs modélisés par des fonctions affines
\end{enumerate}
\vspace{4pt}

\textbf{\normalsize\color{onbo}Chapitre 10 --- Théorème de Thalès}
\begin{enumerate}[label=\arabic*.,itemsep=1pt,topsep=2pt,leftmargin=1.6em]
\item Identifier une configuration de Thalès et écrire les rapports égaux
\item Calculer une longueur inconnue avec le théorème de Thalès
\item Utiliser la réciproque pour prouver un parallélisme
\item Utiliser la réciproque pour prouver que deux droites ne sont pas parallèles
\item Résoudre un problème de configuration « papillon »
\item Calculer l'effet d'un agrandissement/réduction sur un périmètre
\item Calculer l'effet d'un agrandissement/réduction sur une aire
\end{enumerate}
\vspace{4pt}

\textbf{\normalsize\color{onbo}Chapitre 11 --- Théorème de Pythagore}
\begin{enumerate}[label=\arabic*.,itemsep=1pt,topsep=2pt,leftmargin=1.6em]
\item Calculer l'hypoténuse connaissant les deux côtés de l'angle droit
\item Calculer un côté de l'angle droit connaissant l'hypoténuse et l'autre côté
\item Simplifier un résultat sous forme de racine carrée
\item Utiliser la réciproque pour démontrer qu'un triangle est rectangle
\item Utiliser la contraposée pour démontrer qu'un triangle n'est pas rectangle
\item Calculer une distance entre deux points dans un repère
\item Résoudre un problème concret avec le théorème de Pythagore
\end{enumerate}
\vspace{4pt}

\textbf{\normalsize\color{onbo}Chapitre 12 --- Trigonométrie dans le triangle rectangle}
\begin{enumerate}[label=\arabic*.,itemsep=1pt,topsep=2pt,leftmargin=1.6em]
\item Identifier côté opposé, adjacent et hypoténuse selon l'angle considéré
\item Calculer une longueur connaissant un angle et l'hypoténuse
\item Calculer une longueur connaissant un angle et un autre côté
\item Calculer un angle connaissant deux longueurs
\item Utiliser la relation $\cos^{2}\widehat B+\sin^{2}\widehat B=1$
\item Résoudre un problème de hauteur (angle d'élévation)
\item Résoudre un problème de distance inaccessible
\end{enumerate}
\vspace{4pt}

\textbf{\normalsize\color{onbo}Chapitre 13 --- Translations et vecteurs}
\begin{enumerate}[label=\arabic*.,itemsep=1pt,topsep=2pt,leftmargin=1.6em]
\item Reconnaître si deux vecteurs sont égaux
\item Construire l'image d'un point par une translation
\item Calculer les coordonnées d'un vecteur
\item Calculer la norme d'un vecteur
\item Additionner deux vecteurs
\item Utiliser la relation de Chasles pour simplifier une somme
\item Démontrer qu'un quadrilatère est un parallélogramme
\end{enumerate}
\vspace{4pt}

\textbf{\normalsize\color{onbo}Chapitre 14 --- Angle inscrit, angle au centre}
\begin{enumerate}[label=\arabic*.,itemsep=1pt,topsep=2pt,leftmargin=1.6em]
\item Calculer un angle au centre connaissant l'angle inscrit
\item Calculer un angle inscrit connaissant l'angle au centre
\item Utiliser l'égalité des angles inscrits interceptant le même arc
\item Reconnaître un triangle rectangle inscrit dans un cercle
\item Utiliser la réciproque
\item Résoudre un problème combinant plusieurs propriétés
\item Calculer les angles du triangle isocèle $OAB$
\end{enumerate}
\vspace{4pt}

\textbf{\normalsize\color{onbo}Chapitre 15 --- Solides usuels : sections planes et volumes}
\begin{enumerate}[label=\arabic*.,itemsep=1pt,topsep=2pt,leftmargin=1.6em]
\item Calculer l'aire latérale et le volume d'un cylindre
\item Calculer le volume d'une pyramide
\item Calculer l'aire latérale d'une pyramide régulière
\item Calculer la génératrice et les aires d'un cône
\item Calculer l'aire et le volume d'une sphère ou d'une boule
\item Étudier une section parallèle à la base
\item Résoudre un problème concret (réservoir, emballage)
\end{enumerate}
\vspace{4pt}

\textbf{\normalsize\color{onbo}Chapitre 16 --- Statistiques}
\begin{enumerate}[label=\arabic*.,itemsep=1pt,topsep=2pt,leftmargin=1.6em]
\item Construire un tableau d'effectifs et de fréquences à partir d'une liste brute
\item Calculer les effectifs et fréquences cumulés croissants
\item Calculer une moyenne pondérée par des effectifs
\item Calculer une moyenne pondérée par des coefficients
\item Déterminer une médiane
\item Calculer l'étendue d'une série
\item Construire ou lire un diagramme circulaire
\end{enumerate}
\vspace{4pt}

\end{multicols}
